The profile likelihood approach begins by introducing a parameterization based on strength parameters:
\begin{equation}\label{eq:likelihood_parameterization}
\mathcal{L}(\mu_{s}, \mu_{b}) = P(N| \mu_{s} + \mu_{b} \cdot b^{\mathrm{exp}}_{\mathrm{SR}}) \times P(M|\mu_{b} \cdot b_{\mathrm{CR}}^{\mathrm{exp}}),
\end{equation}
where, $b^{\mathrm{exp}}_{\mathrm{SR}}$ and $b^{\mathrm{exp}}_{\mathrm{CR}}$ denote the expected background yields in the signal and control regions, respectively, as determined from MC simulation. The parameter $\mu_{s}$ represents the signal strength, while $\mu_{b}$ represents the background strength.

When incorporating a NP $\theta$, a constraint term $N(\tilde{\theta}|\theta)$ can be introduced to represent an auxiliary measurement $\tilde{\theta}$ associated with $\theta$. In this analysis, two general forms are used for such constraints:
\begin{itemize}
  \item A unit Gaussian: $G(\tilde{\theta}|\theta,1) = \frac{1}{\sqrt{2\pi}} , e^{-\frac{{(\tilde{\theta}-\theta)}^{2}}{2}}$
  \item A Poisson: $P(\tilde{\theta}|\theta \lambda) = \frac{{(\theta \lambda)}^{\tilde{\theta}} e^{-\theta \lambda}}{\tilde{\theta}!}$
\end{itemize}
where $\lambda$ is a constant, typically taken as the nominal value of $\tilde{\theta}$.

Our signal and background expectations are defined as functions of the NPs $\theta$. The functions are parameterized such that the response of s and b to each $\theta$ is factorized from the nominal value of the expected rate. In mathematical terms, this is $s = s_{0} \times \prod \nu(\theta)$ where the form of $\nu(\theta)$ depends on the source of the systematic uncertainty. There are four cases to determine the form of $\nu(\theta)$:
\begin{itemize}
  \item \textbf{Flat systematics} --- Systematics that do not change the shape of the discriminating variable take the form $\nu_{\mathrm{flat}} = \kappa^{\theta}$, where $\kappa$ is determined from $\nu_{\mathrm{float}}$ at $\theta = \pm{1}$. This corresponding likelihood constraint on $\theta$ is a unit Gaussian, with $\kappa^{\theta}$ following a lognormal distribution.
  \item \textbf{Shape-affecting systematics} --- If a systematic affects both the normalization and shape, it is decomposed into a flat component (treated as above) and a pure shape component that does not alter the total rate. The shape variation is modeled using vertical linear interpolation $\nu_{\mathrm{shape}} = 1 + \epsilon\theta$, where $\epsilon$ is measured from $\nu_{shape}$ at $\theta = \pm{1}$. The constraint is a unit Gaussian, and a truncation is applied so that $\nu_{\mathrm{shape}}(\theta) = 0$ for $\theta < -1/\epsilon$. Both components share the same $\theta$ parameter.
  \item \textbf{MC statistical / data-driven uncertainties} --- These are modeled with a Poisson constraint $P(\tilde{\theta}|\theta\lambda)$. For an uncertainty $\sigma_{b}$ on an expected rate $b_{0}$ we have $\tilde{\theta} = \lambda = \frac{b_{0}^{2}}{\sigma_{b}^{2}}$ and $\nu_{\mathrm{stat}} = 0$.
  \item \textbf{High statistic control regions} --- Used to constrain the normalization of a background, these are treated similarly to the point before this. The Poisson constraint is $P(\tilde{\theta}|\theta\lambda) = \frac{{(\theta\lambda)}^{\tilde{\theta}}e^{-\theta\lambda}}{\tilde{\theta}!}$ where $\tilde{\theta}$ is the observed events in the control region. The expected number of events is $\lambda = \mu_{s} + \theta b_{\mathrm{target}} + \sum_{i}^{N_{\mathrm{backgrounds}}-1}b_{i}$ where $b_{\mathrm{target}}$ is the background targeted by the control region and the sum accounts for contamination from both signal and other backgrounds. In the full likelihood, different NPs represent the strengths of nearly all backgrounds, and strength parameters multiply the expected yields wherever those backgrounds appear, ensuring proper treatment of cross-contamination among control regions.
\end{itemize}

Each $\theta$ corresponds to a distinct source of systematic uncertainty and may affect multiple signal and background yields in a correlated way, though not all backgrounds are impacted by every $\theta$. For example, the QCD scale uncertainty on $q\bar{q} \to WZ$ does not affect the top background.

For shape systematics, only statistically significant variations are included. Variations arising from limited MC statistics can lead to convergence issues during minimization, as the shape extrapolation outside the $\pm 1\sigma$ range of the NP can become unstable.
